\begingroup
\renewcommand{\arraystretch}{1.45}
\fancyhead[L]{\small DETERMINATION OF ATTRACTION}
\Needspace{29\baselineskip}
\AttractionArticle{8}
Although formulas 15, 16 could suffice for determining the coefficients $\gamma,\gamma',\gamma''$, still more elegant ones may be given. To this end we shall multiply equation [5] by $aabb-GG$, whence, after a slight reduction, the following results
\[
\begin{aligned}
&\frac{aaAA(bb+G)}{aa+G}-AAG+\frac{bbBB(aa+G)}{bb+G}-BBG+\frac{aabbCC}{G}-CCG\\[6pt]
&\qquad{}=aabb-GG
\end{aligned}
\]
But from the nature of the cubic equation it follows that
\[
\begin{aligned}
&\text{sum of the roots}\quad G-G'-G''=AA+BB+CC-aa-bb\\[8pt]
&\text{product of the roots}\quad GG'G''=aabbCC
\end{aligned}
\]
\Needspace{24\baselineskip}
Hence the preceding equation passes into the following:
\[
\begin{aligned}
&\frac{aaAA(bb+G)}{aa+G}+\frac{bbBB(aa+G)}{bb+G}+G'G''-G(G-G'-G''+aa+bb)\\[6pt]
&\qquad{}=aabb-GG
\end{aligned}
\]
which may also be written as follows
\[
\begin{aligned}
&\frac{aaAA(bb+G)}{aa+G}+\frac{bbBB(aa+G)}{bb+G}-(aa+G)(bb+G)\\[6pt]
&\qquad{}+(G+G')(G+G'')=0
\end{aligned}
\]
\Needspace{30\baselineskip}
Hence the value of the coefficient $\gamma$ from the first formula in [15] is transformed into the following:
\[
\begin{aligned}
&\gamma=\surd{\frac{(aa+G)(bb+G)}{(G+G')(G+G'')}}\quad.\,.\,.\,.\,.\,.\,.\,.\quad[17]
\end{aligned}
\]
By an entirely similar analysis the following is found
\[
\begin{aligned}
&\gamma'=\surd{\frac{(aa-G')(bb-G')}{(G+G')(G''-G')}}\quad.\,.\,.\,.\,.\,.\,.\,.\quad[18]
\end{aligned}
\]
\[
\begin{aligned}
&\gamma''=\surd{\frac{(aa-G'')(bb-G'')}{(G+G'')(G'-G'')}}\quad.\,.\,.\,.\,.\,.\,.\quad[19]
\end{aligned}
\]
\SourceLeaf{354}
\par After the coefficients $\gamma,\gamma',\gamma''$ have been found, the remaining ones $\alpha,\beta,\alpha',\beta',\alpha'',\beta''$ will be derived from them by formulas 3, 4, 6, 7, 9,\nobreak{} 10.

\Needspace{43\baselineskip}
\AttractionArticle{9}
It is readily seen that the signs of the radical expressions by which we have determined $\gamma,\gamma',\gamma''$ may be chosen at will. It is worth investigating, however, how the sign of the quantity $\varepsilon$ is connected with those signs. To this end let us consider the third equation in III, Art. 3.
\[
\begin{aligned}
&\varepsilon\gamma=\alpha'\beta''-\beta'\alpha''
\end{aligned}
\]
which is transformed by formulas 6, 7, 9, 10 into this:
\[
\begin{aligned}
&\varepsilon\gamma=\frac{abAB\gamma'\gamma''}{(aa-G')(bb-G'')}-\frac{abAB\gamma'\gamma''}{(aa-G'')(bb-G')}\\[8pt]
&=\frac{ab(aa-bb)AB(G''-G')\gamma'\gamma''}{(aa-G')(aa-G'')(bb-G')(bb-G'')}
\end{aligned}
\]
But from consideration of equation 13 we readily deduce
\[
\begin{aligned}
&(aa+G)(aa-G')(aa-G'')=aa(aa-bb)AA\\[8pt]
&(bb+G)(bb-G')(bb-G'')=-bb(aa-bb)BB
\end{aligned}
\]
\Needspace{22\baselineskip}
Hence the preceding equation becomes
\[
\begin{aligned}
&\varepsilon\gamma=\frac{(aa+G)(bb+G)(G'-G'')\gamma'\gamma''}{ab(aa-bb)AB}
\end{aligned}
\]
which, combined with equation 17, gives
\[
\begin{aligned}
&\gamma\gamma'\gamma''=\frac{\varepsilon ab(aa-bb)AB}{(G+G')(G+G'')(G'-G'')}
\end{aligned}
\]
\Needspace{33\baselineskip}
Hence it is evident that, if the negative root of the cubic equation that is greater in absolute value has been chosen for $-G'$, and at the same time the coefficients $\gamma,\gamma',\gamma''$ have all been taken positive, $\varepsilon$ acquires the same sign as $AB$; and that the same happens if contrary choices have been made for all four of these conditions or for two of them, but the opposite if you have reversed one or three conditions. The following additional relations, which are easily derived from the preceding ones, should also be noted:
\SourceLeaf{355}
\[
\begin{aligned}
&\alpha\alpha'\alpha''=\frac{\varepsilon aabAAB}{(G+G')(G+G'')(G'-G'')}\\[8pt]
&\beta\beta'\beta''=-\frac{\varepsilon abbABB}{(G+G')(G+G'')(G'-G'')}\\[8pt]
&\alpha\beta=\frac{abAB}{(G+G')(G+G'')}\\[8pt]
&\alpha'\beta'=-\frac{abAB}{(G+G')(G'-G'')}\\[8pt]
&\alpha''\beta''=\frac{abAB}{(G+G'')(G'-G'')}
\end{aligned}
\]

\clearpage
\endgroup
